
\documentclass{article}
\usepackage{colortbl}
\usepackage{makecell}
\usepackage{multirow}
\usepackage{supertabular}

\begin{document}

\newcounter{utterance}

\centering \large Interaction Transcript for game `cladder', experiment `full\_v1.5\_default', episode 8830 with qwen3.5{-}9b.
\vspace{24pt}

{ \footnotesize  \setcounter{utterance}{1}
\setlength{\tabcolsep}{0pt}
\begin{supertabular}{c@{$\;$}|p{.15\linewidth}@{}p{.15\linewidth}p{.15\linewidth}p{.15\linewidth}p{.15\linewidth}p{.15\linewidth}}
    \# & \multicolumn{2}{c}{Player} && \multicolumn{2}{c}{Game Master} \\
    \hline

    \theutterance \stepcounter{utterance}  
    & & & \multicolumn{4}{p{0.6\linewidth}}{
        \cellcolor[rgb]{0.9,0.9,0.9}{
            \makecell[{{p{\linewidth}}}]{
                \texttt{\tiny{[P1$\langle$GM]}}
                \texttt{You are an expert in causal inference. The following question is not a typical commonsense query, but rather a meticulously designed question created by a professor specializing in causal inference, intended to assess the students' mastery of the course content.} \\
\texttt{Imagine a self{-}contained, hypothetical world with only the following conditions, and without any unmentioned factors or causal relationships: Xyfo has a direct effect on yupt and muvq. Zupj has a direct effect on yupt. Yupt has a direct effect on muvq. Xyfo is unobserved. Method 1: We look at how yupt correlates with muvq case by case according to xyfo. Method 2: We look directly at how yupt correlates with muvq in general. To understand how yupt affects muvq, is it more correct to use the Method 1 than Method 2?} \\
\texttt{Start your answer with 'Yes' or 'No', followed by additional reasoning or evidence to support your explanation.} \\
            }
        }
    }
     \\ \\

    \theutterance \stepcounter{utterance}  
    & \multicolumn{4}{p{0.6\linewidth}}{
        \cellcolor[rgb]{0.561,0.737,0.561}{
            \makecell[{{p{\linewidth}}}]{
                \texttt{\tiny{[P1$\rangle$GM]}}
                \texttt{No} \\
\\ 
\texttt{**Reasoning:**} \\
\\ 
\texttt{In this hypothetical world, the causal structure can be represented as follows:} \\
\texttt{1.  \$X\_\{yfo\} \textbackslash{}rightarrow Y\_\{upt\}\$} \\
\texttt{2.  \$X\_\{yfo\} \textbackslash{}rightarrow M\_\{uvq\}\$} \\
\texttt{3.  \$Z\_\{upj\} \textbackslash{}rightarrow Y\_\{upt\}\$} \\
\texttt{4.  \$Y\_\{upt\} \textbackslash{}rightarrow M\_\{uvq\}\$} \\
\\ 
\texttt{The goal is to estimate the causal effect of \$Y\_\{upt\}\$ on \$M\_\{uvq\}\$.} \\
\\ 
\texttt{**Analysis of Method 2 (General Correlation):**} \\
\texttt{Method 2 looks at the correlation between \$Y\_\{upt\}\$ and \$M\_\{uvq\}\$ without conditioning on \$X\_\{yfo\}\$. Because \$X\_\{yfo\}\$ is an unobserved common cause (a confounder) of both \$Y\_\{upt\}\$ and \$M\_\{uvq\}\$, it induces a spurious association between them. The observed correlation in Method 2 will be a mixture of the true causal effect of \$Y\_\{upt\}\$ on \$M\_\{uvq\}\$ and the bias introduced by the unobserved confounder \$X\_\{yfo\}\$. Therefore, Method 2 yields a biased estimate of the causal effect.} \\
\\ 
\texttt{**Analysis of Method 1 (Stratification/Conditioning):**} \\
\texttt{Method 1 examines the correlation between \$Y\_\{upt\}\$ and \$M\_\{uvq\}\$ "case by case according to \$X\_\{yfo\}\$." This is equivalent to stratifying the data by the levels of \$X\_\{yfo\}\$ (or conditioning on \$X\_\{yfo\}\$). By holding \$X\_\{yfo\}\$ constant, we block the back{-}door path \$Y\_\{upt\} \textbackslash{}leftarrow X\_\{yfo\} \textbackslash{}rightarrow M\_\{uvq\}\$. Since \$X\_\{yfo\}\$ is the only unobserved confounder mentioned, and we are conditioning on it (even though it is unobserved in the general population, the method implies we are analyzing subgroups defined by its values), this isolates the direct causal pathway \$Y\_\{upt\} \textbackslash{}rightarrow M\_\{uvq\}\$.} \\
\\ 
\texttt{**Conclusion:**} \\
\texttt{Method 1 correctly addresses the confounding bias caused by the unobserved variable \$X\_\{yfo\}\$ by stratifying on it, whereas Method 2 fails to do so. Therefore, to understand the true causal effect of \$Y\_\{upt\}\$ on \$M\_\{uvq\}\$, Method 1 is indeed more correct than Method 2.} \\
            }
        }
    }
    & & \\ \\

    \theutterance \stepcounter{utterance}  
    & & & \multicolumn{2}{p{0.3\linewidth}}{
        \cellcolor[rgb]{0.9,0.9,0.9}{
            \makecell[{{p{\linewidth}}}]{
                \texttt{\tiny{[GM$|$GM]}}
                \texttt{No} \\
            }
        }
    }
    & & \\ \\

    \theutterance \stepcounter{utterance}  
    & & & \multicolumn{2}{p{0.3\linewidth}}{
        \cellcolor[rgb]{0.9,0.9,0.9}{
            \makecell[{{p{\linewidth}}}]{
                \texttt{\tiny{[GM$|$GM]}}
                \texttt{yes} \\
            }
        }
    }
    & & \\ \\

    \theutterance \stepcounter{utterance}  
    & & & \multicolumn{2}{p{0.3\linewidth}}{
        \cellcolor[rgb]{0.9,0.9,0.9}{
            \makecell[{{p{\linewidth}}}]{
                \texttt{\tiny{[GM$|$GM]}}
                \texttt{game\_result = LOSE} \\
            }
        }
    }
    & & \\ \\

\end{supertabular}
}

\end{document}
